% Part: conditional-logics
% Chapter: minimal-change-semantics
% Section: sphere-models
%READERNOTE{SOL6-A214: रिकाम्या गोलकाची परवानगी व अनंत उपकुलांवरील closure स्पष्ट केली. सर्वांत लहान पूर्वांग-गोलक नसण्याचा पूर्ण प्रतिमान दिला; फक्त उतरत्या क्रमिकेवरून तो निष्कर्ष काढता येत नाही. पूर्वीचा लहान पूर्वांग-गोलकांसाठीचा source repair जपला.}

\documentclass[../../../include/open-logic-section]{subfiles}

\begin{document}

\olfileid{con}{min}{sph}

\olsection{गोलक प्रतिमाने}

प्रतिवास्तविक विधानांना औपचारिक
चिन्हार्थमीमांसा देण्याचा एक मार्ग
म्हणजे लुईस यांच्या अनौपचारिक
मांडणीचे गणिती रचनेत रूपांतर करणे.
मग जग~$w$ भोवतीचे गोलक
हे जगांचे संच असतात.
गोलक एकमेकांच्या आत
असल्यामुळे, जग~$w$ भोवतीच्या
या जगांच्या संचांना उपसंच-संबंधाने
रेषीय क्रमात असावे लागते.

\begin{defn}
  \emph{गोलक प्रतिमान} म्हणजे
  $\mModel{M} = \tuple{W, O, V}$
  हे त्रिक: येथे $W$ हा जगांचा
  रिक्त नसलेला संच आहे,
  $V\colon \PVar \to \Pow{W}$
  हे सत्य-मूल्यांकन आहे, आणि
  $O\colon W \to \Pow{\Pow{W}}$
  प्रत्येक जग~$w$ ला
  \emph{गोलकांची प्रणाली}~$O_w$
  नेमून देते. प्रत्येक~$w$ साठी
  $O_w$ हा जगांच्या संचांचा संच
  असून त्याने पुढील अटी पूर्ण
  केल्या पाहिजेत:
  \begin{enumerate}
  \item $O_w$ चे \emph{केंद्र}~$w$
    असते: $\{w\} \in O_w$.
  \item $O_w$ मधील गोलक
    \emph{एकमेकांच्या आत} असतात:
    $S_1$, $S_2 \in O_w$ असतील,
    तर $S_1 \subseteq S_2$ किंवा
    $S_2 \subseteq S_1$; म्हणजे
    $O_w$ ला~$\subseteq$ ने
    रेषीय क्रम असतो.
  \item प्रत्येक रिक्त नसलेल्या उपकुल $\mathcal S\subseteq O_w$
    साठी $\bigcup\mathcal S\in O_w$.
  \item प्रत्येक रिक्त नसलेल्या उपकुल $\mathcal S\subseteq O_w$
    साठी $\bigcap\mathcal S\in O_w$.
    या अटी अनंत उपकुलांनाही लागू होतात.
  \end{enumerate}
\end{defn}

$O_w$ मागील कल्पना अशी:
$w$ भोवतीची जगे~$w$ पासून
किती दूर आहेत यानुसार
थरांत मांडली जातात.
रिकामे गोलक वगळल्यास सर्वांत आतला गोलक म्हणजे
एकटे~$w$, अर्थात~$\{w\}$: इतर कोणत्याही जगापेक्षा
$w$ हे स्वतःच्या अधिक जवळ असते. व्याख्या रिकामा
गोलकही परवानगी देते; त्याने पुढील सत्यता-अट बदलत नाही.
$S \subsetneq S'$ असेल,
तर $S' \setminus S$
मधील जगे~$S$ मधील
जगांपेक्षा~$w$ पासून
अधिक दूर असतात:
$S'\setminus S$ हा $S'$ च्या आत पण $S$ च्या बाहेर
असलेल्या जगांचा ‘‘थर’’ आहे.
विशेषतः, गोलकांत
त्यांच्या बाह्य पृष्ठभागाच्या
आतील सर्व जगे असतात
असे समजावे; गोलक
म्हणजे फक्त स्वतंत्र
थर नव्हेत.

\begin{figure}
\begin{center}
\begin{tikzpicture}[layerwidth=1.5,scale=.6]\tiny
  \spheresystem{3}
  \propositionintersect{0}{3}{60}{6}
  \path (0,0) node[world] {$w$};
  \spherepos{-30}{2}{node[world] {$w_2$}}
  \spherepos{220}{2}{node[world] {$w_3$}}
  \spherepos{90}{2}{node[world] {$w_1$}}
  \spherepos[shift={(.1,0)}]{10}{3}{node[world] {$w_5$}}
  \spherepos[shift={(.1,0)}]{-10}{3}{node[world] {$w_6$}}
  \spherepos{225}{3}{node[world] {$w_4$}}
  \spherepos{20}{4}{node[world] {$w_7$}}
  \path (5,0) node[left] {$p$};
\end{tikzpicture}
\caption{गोलक प्रतिमानाची आकृती}
\ollabel{fig:sphere-model}
\end{center}
\end{figure}

\olref{fig:sphere-model} मधील
आकृती पुढील गोलक प्रतिमान
दाखवते: $W = \{w, w_1, \dots, w_7\}$,
$V(p) = \{w_5, w_6,
w_7\}$. सर्वांत आतला गोलक
$S_1 = \{w\}$ आहे.
$w$ च्या सर्वाधिक जवळची
जगे $w_1, w_2, w_3$
आहेत; म्हणून पुढचा मोठा
गोलक $S_2 = \{w, w_1,
w_2, w_3\}$ आहे.
त्याहून बाहेरची जगे
$w_4$, $w_5$, $w_6$
आहेत; म्हणून सर्वांत
बाहेरचा गोलक
$S_3 = \{w, w_1, \dots, w_6\}$
आहे. $w$ भोवतीची
गोलक-प्रणाली
$O_w = \{S_1, S_2, S_3\}$
आहे. जग~$w_7$ हे~$w$
भोवतीच्या कोणत्याही गोलकात
नाही. ज्या जगांत~$p$
सत्य आहे त्यांपैकी
सर्वाधिक जवळची जगे
$w_5$ आणि~$w_6$
आहेत. त्यामुळे~$p$
सत्य असलेले जग
समाविष्ट करणारा सर्वांत
लहान गोलक~$S_3$ आहे.

गोलक प्रतिमान~$\mModel M$
मधील जग~$w$ येथे
सूत्र~$!A$ ची पूर्ति,
$\mSat{M}{!A}[w]$,
याची व्याख्या करताना
मोडल !!{formula}s
यांच्या व्याख्येत
$!B \cif !C$
साठी पुढील अट जोडतो:

\begin{defn}
  $\mSat{M}{!B \cif !C}[w]$
  तेव्हा आणि केवळ तेव्हाच,
  जेव्हा पुढीलपैकी एक घडते:
  \begin{enumerate}
  \item\ollabel{sphere-vac}
    सर्व $u \in \bigcup O_w$
    साठी $\mSat/{M}{!B}[u]$; किंवा
  \item\ollabel{sphere-nonvac}
    एखाद्या $S \in O_w$ साठी,
    \begin{enumerate}
      \item एखाद्या $u \in
        S$ साठी
        $\mSat{M}{!B}[u]$; आणि
      \item सर्व $v \in S$
        साठी, एकतर
        $\mSat/{M}{!B}[v]$
        किंवा
        $\mSat{M}{!C}[v]$.
    \end{enumerate}
  \end{enumerate}
\end{defn}

या व्याख्येनुसार
$\mSat{M}{!B \cif !C}[w]$
तेव्हा आणि केवळ तेव्हाच, जेव्हा
एकतर पूर्वांग~$!B$
हे~$w$ भोवतीच्या सर्व
गोलकांत सर्वत्र असत्य असते,
किंवा एखाद्या गोलक~$S$
मध्ये~$!B$ सत्य असते आणि
त्या ‘‘$!B$ समाविष्ट
करणाऱ्या’’ गोलकातील
सर्व जगांत
$!B \lif !C$
हे वास्तविक अभिव्यंजन
सत्य असते.
सहज स्पष्टीकरणावरून
अपेक्षा वाटू शकते,
तरी व्याख्येत~$S$
हा~$!B$ समाविष्ट
करणारा \emph{सर्वांत
आतला} गोलक असावा
अशी अट घातलेली नाही.
पण \olref{sphere-nonvac}
मधील अट एखाद्या
गोलक~$S$ साठी
पूर्ण होत असेल,
तर~$S$ मध्ये
समाविष्ट असलेल्या आणि
पूर्वांग सत्य असलेले जग
धारण करणाऱ्या प्रत्येक
लहान गोलकासाठीही
ती पूर्ण होते.
म्हणून असा सर्वांत
आतला गोलक अस्तित्वात
असेल, तर त्यासाठीही
अट पूर्ण होते.

गोलक प्रतिमानाच्या व्याख्येनुसार $!B$ सत्य असलेले
जग धारण करणारा सर्वांत आतला गोलक असतोच असे नाही.
केवळ अनंत उतरती क्रमिका दिल्याने तो नसल्याचे सिद्ध होत नाही;
प्रणालीत तिच्यापेक्षा लहान पूर्वांग-गोलकही असू शकतो.
एक पूर्ण उदाहरण असे: $W=\{w\}\cup\{u_n:n\geq1\}$,
$S_n=\{w\}\cup\{u_j:j\geq n\}$ आणि
$O_w=\{\{w\}\}\cup\{S_n:n\geq1\}$ घ्या.
इतर जगांसाठी $O_{u_n}=\{\{u_n\}\}$ घ्या आणि
पूर्वांग म्हणून $!B=p$, $V(p)=\{u_n:n\geq1\}$ ठेवा.
ही प्रणाली केंद्रित, एकमेकांच्या आत असलेली आणि प्रत्येक
रिक्त नसलेल्या उपकुलाच्या संघटनाखाली व छेदाखाली बंद आहे.
गोलक $S_n$ हे सर्व $p$-जग धारण करतात,
$S_{n+1}\subsetneq S_n$ आणि $\bigcap_{n\geq1}S_n=\{w\}$,
पण $p$ हे $w$ येथे असत्य आहे. म्हणून सर्वांत लहान
$p$-जग धारण करणारा गोलक नाही.
या प्रतिमानात $\mSat{M}{p\cif !C}[w]$ तेव्हा आणि केवळ तेव्हाच,
जेव्हा एखाद्या $i\geq1$ साठी सर्व $j\geq i$ येथे
$\mSat{M}{!C}[u_j]$. सममूल्यपणे, अशा $i$ पासून
$S_i,S_{i+1},\dots$ या सर्व गोलकांत $p\lif !C$ सत्य असते.
म्हणून सर्वांत लहान पूर्वांग-गोलक नसतानाही मूळ सत्यता-अट
निर्णायक आहे.

\end{document}
